\section{Cavity Count, Volume, and Thermal Compliance}
\label{app:cavitybudget}

The following facts explain the budget experiment; they are properties of the
classical ideal Dirichlet cooling model, not new topological-derivative claims.

\begin{proposition}[Monotonicity under additional cooled material removal]
\label{prop:coolingmonotonicity}
Let $\Omega\subset\mathbb R^3$ be bounded and Lipschitz, let $f\in L^2(\Omega)$,
and let $D_2\subset D_1\subset\Omega$ be nonempty Lipschitz material domains.
For $i=1,2$, let $u_i\in H_0^1(D_i)$ solve
$-\Delta u_i=f|_{D_i}$ weakly, with zero trace on every boundary component.
Extend $u_i$ by zero to $\Omega$ and define
$J_i=\int_\Omega f u_i$. Then
\begin{equation}
 J_1-J_2=\int_\Omega |\nabla(u_1-u_2)|^2\ge0.
 \label{eq:coolingprojection}
\end{equation}
Equality holds if and only if the extended states agree. If $D_1\setminus D_2$
contains a nonempty open set on which $f>0$ almost everywhere, the inequality
is strict. Division by the fixed volume $|\Omega|$ preserves these conclusions.
\end{proposition}
\begin{proof}
\textbf{Step 1: common function space.}
By the definition of $H_0^1$, each $u_i$ is an $H^1$ limit of smooth functions
compactly supported in $D_i$. Extending these approximations by zero proves
that their extensions belong to $H_0^1(\Omega)$. The same argument gives the
inclusion of extended spaces $H_0^1(D_2)\subset H_0^1(D_1)$. Their gradients
vanish almost everywhere outside their respective material domains.

\textbf{Step 2: existence and energy identities.}
The bilinear form $a(v,w)=\int_\Omega\nabla v\cdot\nabla w$ is coercive on
each extended space by Poincare's inequality on $\Omega$; the right-hand side
$\ell(v)=\int_\Omega fv$ is bounded by Cauchy--Schwarz and Poincare.
Lax--Milgram gives a unique weak solution in each space. Testing its weak
equation with $u_i$ gives $a(u_i,u_i)=\ell(u_i)=J_i$.

\textbf{Step 3: the cross term.}
The space inclusion makes $u_2$ an admissible test function in the equation
for $u_1$. Hence $a(u_1,u_2)=\ell(u_2)=J_2$. Expanding the squared gradient norm
therefore gives $a(u_1-u_2,u_1-u_2)=J_1-2J_2+J_2$, proving
\cref{eq:coolingprojection}. A zero norm forces $u_1=u_2$ by Poincare; the
converse is immediate.

\textbf{Step 4: strictness under positive interior heating.}
If equality held and $U\subset D_1\setminus D_2$ were the stated open set,
then $u_1=u_2=0$ almost everywhere on $U$, so $\nabla u_1=0$ there. Choose a
nonnegative, nonzero $\phi\in C_c^\infty(U)$. The first weak equation would
give $0=\int_U f\phi>0$, a contradiction. Finally, multiplication by the
positive constant $|\Omega|^{-1}$ proves the normalized version.
\end{proof}

This proposition applies to nested removals, such as increasing the radius
of fixed centers or adding a new cavity while retaining old ones. Replacing
two large cavities by twelve small ones at the same removed volume does not
produce nested material domains; monotonicity alone cannot rank those designs.
Insulating, Robin, or fluid-coupled cavity walls require a different analysis.

\begin{proposition}[Exact geometry at fixed removed volume]
\label{prop:budgetgeometry}
Let $n$ disjoint spherical cavities of equal radius $r_n$ have total volume
$V>0$. Their radius, total exposed area $A_n$, and sum $S_n$ of isolated
Newtonian capacities (energy convention) are
\begin{align}
 r_n&=\left(\frac{3V}{4\pi n}\right)^{1/3},\\
 A_n&=(36\pi)^{1/3}V^{2/3}n^{1/3},\\
 S_n&=4\pi\left(\frac{3V}{4\pi}\right)^{1/3}n^{2/3}.
 \label{eq:budgetgeometry}
\end{align}
Here $S_n$ is not the capacity of the union and does not give its exact cooling
benefit.
\end{proposition}
\begin{proof}
Disjointness makes the removed volume the sum of individual ball volumes:
$V=n(4\pi/3)r_n^3$. Solving gives the first formula. Each sphere has area
$4\pi r_n^2$, and separation leaves that entire surface exposed; substitution
into $n4\pi r_n^2$ gives the second formula. For completeness, the equilibrium
potential of a radius-$r$ ball in $\mathbb R^3$ is $1$ inside and $r/|x|$
outside, after translation to the origin. Its energy is
\[
 \int_{\mathbb R^3}|\nabla P|^2
 =4\pi\int_r^\infty \frac{r^2}{s^4}s^2\,\mathrm ds=4\pi r.
\]
To check the minimizing property, write any smooth exterior competitor with
the same trace and decay as $P+w$, where $w$ has zero trace at $|x|=r$.
Integration by parts gives $\int\nabla P\cdot\nabla w=0$, since $P$ is
harmonic outside the ball and the boundary terms vanish. Thus its energy is
$\int|\nabla P|^2+\int|\nabla w|^2\ge\int|\nabla P|^2$; density extends the
identity to finite-energy, zero-trace perturbations. For the equivalent
capacity definition requiring at least one on the ball, truncating a
competitor at one cannot increase energy and reduces to this boundary-value
problem.
Summing these individual energies and substituting $r_n$ proves the third
formula. This sum treats every ball in isolation; their joint equilibrium
potential satisfies a different exterior boundary-value problem, so no union
capacity identity follows.
\end{proof}

Thus equal removed volume does not equalize cooling area. Classical
small-cavity scaling suggests why fragmentation can help, but the
single-cavity expansion is not asserted uniformly as $n\to\infty$ with
shrinking separations. Field depletion, interactions, and geometric feasibility
must be evaluated at the actual finite configurations. In particular, a
penalty $\lambda V$ is constant across a fixed-volume count sweep and cannot
resolve its count preference. A cavity-count cost, surface-area cost,
minimum radius, wall constraint, or explicit cooling-system model may do so;
this experiment does not identify physical values for those costs.

\paragraph{A finite-design cost envelope, not a continuous optimum.}
For the finite tested designs $d$, let $q_d=J_d/J_0$ and let $n_d$ be the
cavity count. Given a stipulated normalized per-cavity cost $\kappa\ge0$,
minimize $q_d+\kappa n_d$ over this finite set. Designs $a,b$ with
$n_a\ne n_b$ tie precisely at
$\kappa=(q_a-q_b)/(n_b-n_a)$, obtained by rearranging the two affine costs.
The lower envelope can change minimizer only at such pairwise intersections.
Indeed, away from intersections all pairwise differences keep their signs,
so their ordering is constant on each open interval. This construction
describes sensitivity to an assumed cost; neither sampling uncertainty nor
the untested geometry space disappears when an envelope is plotted.
